\section{Operator-Unreadable Remote Relay}
\label{sec:pilot-mobile-opaque-relay}

Pilot completes \texttt{MOB-0308} with an optional store-and-forward relay whose
operator transports opaque envelopes but does not receive the application plaintext,
the mother's private transport key, phone signing key, capability lease authority or
provider credentials.  The relay is a continuity path, not a second mother brain.

\subsection{Route establishment}

When an owner configures \texttt{PILOT\_RELAY\_ENDPOINT}, the mother creates a persistent
X25519 transport key and an opaque random route identifier.  Both are stored in the
customer-controlled runtime with owner-only file permissions.  The private transport key
never leaves the mother.  Successful local phone pairing returns a short-lived,
mother-signed \texttt{pilot.opaque-relay-route.v1} descriptor containing:
\begin{itemize}
  \item the relay HTTPS endpoint and an opaque route identifier;
  \item the mother's X25519 public transport key;
  \item a bounded submission capability, expiry and request-size limits;
  \item the fixed encryption suite and the mother's identity fingerprint; and
  \item explicit declarations that the operator cannot read content and that no mother
        secret is embedded.
\end{itemize}

Changing the endpoint, transport key, route identifier, bounds or security declarations
invalidates the mother's Ed25519 signature.  A route is accepted by the phone only when
it matches the mother identity established during physical or trusted-LAN pairing.

\subsection{Encrypted request and response}

The phone generates a fresh ephemeral X25519 key for every relayed operation.  It derives
a shared secret with the signed mother public transport key.  HKDF-SHA-256 produces
separate AES-256-GCM keys for the request and response directions.  The requested API
path, phone certificate, lease, possession signature and private body are all inside the
encrypted payload.  Only the following routing metadata remains visible:
\begin{itemize}
  \item opaque route and request identifiers;
  \item creation and expiry times;
  \item the ephemeral public key;
  \item bounded ciphertext length; and
  \item unavoidable connection metadata such as IP address and timing.
\end{itemize}

For shared secret $S$, request identifier $q$, direction $d$ and authenticated metadata
$A$, the payload is
\[
 K_d=\operatorname{HKDF}_{\mathrm{SHA256}}(S,\,\texttt{PilotRelay}\parallel q\parallel d),
 \qquad C_d=\operatorname{AESGCM}_{K_d}(N_d,P_d,A).
\]
The mother opens the request, applies the normal certificate, lease, scope, revocation,
signature, approval and idempotency checks, then seals the exact status and response body
to the originating ephemeral key.  The relay cannot forge a successful response because
it cannot calculate the response-direction key or authentication tag.

\subsection{Fallback and failure semantics}

The phone attempts the trusted LAN first, then each signed customer-owned direct endpoint.
The opaque relay is attempted only after those routes fail at the network layer.  A policy
denial, expired lease, invalid signature or ordinary application error is authoritative and
never triggers route hopping.  Consequential requests retain their original idempotency
keys, so uncertain transport does not authorize duplicate actions.

Requests expire after at most two minutes.  The reference queue bounds pending envelopes,
deduplicates request identifiers, rejects invalid capabilities and deletes a response when
the phone receives it.  The relay has no API that returns decrypted content.  A relay can
still observe traffic volume or refuse delivery; it cannot be made trusted for availability
merely by making it blind to content.

\subsection{Verification and deployment boundary}

Automated verification covers signed descriptors, endpoint and size bounds, owner-only key
storage, descriptor tampering, request encryption, response encryption, direction-separated
keys, duplicate submission, expiry, wrong relay capabilities and the phone's strict
LAN--direct--relay order.  The optimized mobile application compiles with the complete
WebCrypto X25519, HKDF and AES-GCM path.

All 564 AION Fabric tests pass, the optimized mobile production build succeeds, and this
append-only section compiles without warnings.

The protocol and reference bounded queue are implemented, but no Tessaris-operated relay is
silently introduced.  A household or organization must configure an approved HTTPS relay
deployment before this route becomes active.  Automatic network-change detection, bounded
reconnection and live-subscription cursor resumption across Wi-Fi and mobile data remain the
separate \texttt{MOB-0309} stage.
